\documentclass[10pt,a4paper]{amsart}
\usepackage[margin=19mm]{geometry}
\usepackage{amsmath,amssymb,mathtools}
\usepackage[scr=rsfs,cal=euler]{mathalfa}
\usepackage{xfrac}
\usepackage[unicode,hidelinks,bookmarks=false]{hyperref}
\newcommand{\ie}{i.e.\ }
\newcommand{\cf}{cf.\ }
\newcommand{\resp}{respectively\ }
\newcommand{\Subsection}{\subsection}
\providecommand{\nobreakdash}{\nobreak\hbox{-}}
\setlength{\emergencystretch}{2em}
\multlinegap0pt
\usepackage{float}
\usepackage{url}
\newtheorem{theorem}{Theorem}[section]
\newtheorem{prop}[theorem]{Proposition}
\newtheorem{lemma}[theorem]{Lemma}
\newtheorem{corollary}[theorem]{Corollary}
\newtheorem{definition}[theorem]{Definition}
\newtheorem{example}[theorem]{Example}
\newtheorem{remark}[theorem]{Remark}
\usepackage{mathtools}
\usepackage{enumitem}
\usepackage{euscript}
\begin{document}
\title[Closure operators on semilattice-ordered semigroups and spec]{Closure operators on semilattice-ordered semigroups and spectrality of induced operations}
\author{Damian Siejwa}
\date{}
\maketitle
\noindent\emph{Source and licence.} Closure operators on semilattice-ordered semigroups and spectrality of induced operations. Version arXiv:2607.25674v2 (CC BY 4.0). This material is distributed under the Creative Commons Attribution 4.0 International licence, http://creativecommons.org/licenses/by/4.0/, as recorded in the arXiv OAI licence metadata for this exact version and shown on the abstract page https://arxiv.org/abs/2607.25674v2. Full rights notice: licenses/arxiv-2607.25674-CC-BY-4.0.txt.\par\smallskip
\noindent\textit{Editorial scope.} Counted unit: the original \texttt{theorem} environment \texttt{construction\_of\_spectral\_space} at source lines 188--191 together with its complete proof at lines 192--213; the delivered document keeps source lines 134--247 so that every referenced label and every citation resolves inside it. Native editable source is the paper's own concatenated TeX; the AMS wrapper is editorial formatting only. No publisher class, upstream script or unrelated artwork is redistributed.
\par\smallskip\noindent\textit{Reference note.} This article ships no compiled bibliography (biblatex); the entries delivered here are composed from the author's own BibTeX (.bib) fields, with no field invented and no entry dropped.
\begin{theorem}\cite[Corollary 3.3]{Finocchiaro03042014}\label{ultrafilter_criterion}
    A topological space $X$ is spectral if and only if it satisfies the $T_0$-axiom and there is a subbasis $\mathbb{S}$ of $X$ such that 
    \begin{equation*}
        X_{\mathbb{S}}(\mathscr{U}) = \{x \in X \mid \forall S\in\mathbb{S} \hspace{3pt} x\in S \iff S \in\mathscr{U}\} \neq \varnothing
    \end{equation*}
    for any ultrafilter $\mathscr{U}$ on $X$. 
\end{theorem}

Let us recall that given a set $S$, the power set $\mathcal{P}(S)$ endowed with the \emph{hull-kernel topology} is a spectral space. For an arbitrary subset $E \subseteq S$, define
\begin{equation*}
    D(E) := \{A \in \mathcal{P}(S) \mid E \nsubseteq A\}, \qquad V(E) := \{A \in \mathcal{P}(S) \mid E \subseteq A\}. 
\end{equation*}
Thus, $V(E) = \mathcal{P}(S) \setminus D(E)$. The family $\{D(F)\}_{F\in\mathcal{P}_\mathrm{fin}(S)}$ forms an open subbasis of the hull-kernel topology on $\mathcal{P}(S)$. Moreover, each set of this family is quasi-compact in $\mathcal{P}(S)$. For a singleton $\{x\}$ we write $D(x)$ and $V(x)$ instead of $D(\{x\})$ and $V(\{x\})$. 

The \emph{specialization preorder} on a topological space $X$ is defined by 
\begin{equation*}
    A \leqslant B \iff B \in \overline{\{A\}}. 
\end{equation*}
It is well known that in the hull-kernel topology on $\mathcal{P}(S)$ this preorder coincides with set-theoretic inclusion. Indeed, let $A, B \in \mathcal{P}(S)$ and observe that 
\begin{equation*}\begin{aligned}
    B \in \overline{\{A\}} &\iff \forall F\in\mathcal{P}_\mathrm{fin}(S) \hspace{3pt} (F \subseteq A \implies B \in V(F)) \\ &\iff \forall F\in\mathcal{P}_\mathrm{fin}(S) \hspace{3pt} (F \subseteq A \implies F \subseteq B) \\ &\iff A \subseteq B. 
\end{aligned}\end{equation*}
Consequently, the specialization preorder on every subspace $X \subseteq\mathcal{P}(S)$ considered in the rest of the article is set-theoretic inclusion. 
\section{Spectral spaces arising from algebraic closure operators on semilattice-ordered semigroups}\label{spectral_spaces_arising_from_closures}
We begin by considering the family of all subsets that are closed with respect to a given closure operator $\mathrm{cl}$ on a semilattice-ordered semigroup $S$. Namely, we set $$X := \{A\in\mathcal{P}(S) \mid A^\mathrm{cl} = A\}$$ and endow it with the subspace topology induced by the hull-kernel topology on $\mathcal{P}(S)$. For every subset $E \subseteq S$, we write $$D_X(E) := D(E) \cap X, \qquad V_X(E) := V(E) \cap X.$$

Let us recall that quasi-compactness is not in general a hereditary property: if $X_1$ is quasi-compact and $X_2 \subseteq X_1$, then $X_2$ need not be quasi-compact. For this reason, even though each $D(F)$ is quasi-compact in $\mathcal{P}(S)$, it does not follow that the set $D_X(F)$ is quasi-compact in the induced topology on $X$. Lemma \ref{quasi_compactness_of_subbasic_opens} establishes that this property holds whenever the closure operator is algebraic.  
\begin{lemma}\label{quasi_compactness_of_subbasic_opens}
    Let $S$ be a semilattice-ordered semigroup and let $\mathrm{cl}$ be an algebraic closure operator on $S$. For every finite subset $F \subseteq S$, the set $D_X(F)$ is quasi-compact in $X$. 
\end{lemma}
\begin{proof}
    Fix a finite subset $F \subseteq S$. Since $X$ carries the subspace topology induced by the hull-kernel topology on $\mathcal{P}(S)$, the family of sets $$\{D_X(G)\}_{G\in\mathcal{P}_\mathrm{fin}(S)}$$ is an open subbasis of $X$. By Alexander's subbasis theorem, it is enough to show that every cover of $D_X(F)$ by subbasic open sets admits a finite subcover. Suppose then that we have such a cover, i.e.
    \begin{equation}\label{cover_by_subbasic_opens}
        D_X(F) \subseteq \bigcup_{i\in I} D_X(F_i),  
    \end{equation}
    where $F_i$ is a finite subset of $S$ for any $i\in I$. 
    Assume that no finite subcover exists. Then for every finite set $J \subseteq I$, $D_X(F) \nsubseteq \bigcup_{j \in J}D_X(F_j)$ and hence 
    \begin{equation*}
        D_X(F) \setminus\bigcup_{j \in J} D_X(F_j) = D_X(F) \cap \bigcap_{j \in J} V_X(F_j) \neq \varnothing. 
    \end{equation*}
    So for each finite set $J \subseteq I$, we can choose an element $C_J \in X$ such that $F \nsubseteq C_J$ and $\bigcup_{j \in J} F_j \subseteq C_J$. Since $C_J$ is $\mathrm{cl}$-closed and contains $\bigcup_{j\in J}F_j$, it must also contain $(\bigcup_{j \in J} F_j)^{\mathrm{cl}}$. Therefore, if $F \subseteq (\bigcup_{j \in J} F_j)^{\mathrm{cl}}$, then $F\subseteq C_J$, contradicting $F \nsubseteq C_J$. So for every finite set $J \subseteq I$ we have 
    \begin{equation}\label{F_is_not_contained_in_cl_of_finite_sum} 
    F \nsubseteq \Big(\bigcup_{j \in J} F_j\Big)^{\mathrm{cl}}.
    \end{equation}
    
    Now consider the set $A := \bigcup_{i \in I} F_i$. If $F \subseteq A^\mathrm{cl}$, then for each $f \in F$ there exists some finite set $B_f \subseteq A$ such that $f \in B_f^\mathrm{cl}$ (by algebraicity of $\mathrm{cl}$). Let $B := \bigcup_{f\in F}B_f \subseteq A$. Then the set $B$ is finite and hence there exists a finite set $K \subseteq I$ such that $B \subseteq \bigcup_{k \in K} F_k$. We get $$F \subseteq B^\mathrm{cl} \subseteq \Big(\bigcup_{k \in K} F_k\Big)^\mathrm{cl},$$ 
    which leads to a contradiction with \eqref{F_is_not_contained_in_cl_of_finite_sum}. Thus, $F \nsubseteq A^\mathrm{cl}$ and $A^\mathrm{cl} \in D_X(F)$. It is easy to see that also $A^\mathrm{cl} \in \bigcap_{i \in I} V_X(F_i)$. Therefore, 
    \begin{equation*}
        A^\mathrm{cl} \in D_X(F) \cap \bigcap_{i \in I} V_X(F_i) = D_X(F) \setminus \bigcup_{i \in I} D_X(F_i). 
    \end{equation*}
    This is a contradiction with \eqref{cover_by_subbasic_opens}. Thus, every cover of $D_X(F)$ by subbasic open sets must have a finite subcover. By Alexander's subbasis theorem, $D_X(F)$ is quasi-compact. 
\end{proof}
The next theorem is motivated by several earlier spectral space constructions arising from algebraic closure operators. Jun, Ray and Tolliver proved that for an algebraic closure operator $\mathrm{cl}$ the set of $\mathrm{cl}$-closed ideals of a semiring forms a spectral space \cite[Proposition~3.9]{JunRayTolliver}. Related results were obtained for the space of $\mathrm{cl}$-closed submodules of a given module over a commutative ring \cite[Proposition 3.4]{FinocchiaroFontanaSpiritoNullstellensatz}, as well as for the space of stable semistar operations of finite type on an integral domain \cite[Theorem 4.6]{FinocchiaroFontanaSpirito}. Banerjee obtained analogous results in the setting of abelian categories, where he constructed spectral spaces of invariants of closure operators acting on subobjects of a given object \cite[Proposition 3.3]{Banerjee}. Our argument follows the same general philosophy in the context of semilattice-ordered semigroups. In particular, the implication from algebraicity of $\mathrm{cl}$ to spectrality of the associated space $X$ is standard and we include it for completeness. 


\begin{theorem}\label{construction_of_spectral_space}
    Let $S$ be a semilattice-ordered semigroup and let $\mathrm{cl}$ be a closure operator on $S$. Then the space $X = \{A \in\mathcal{P}(S) \mid A^\mathrm{cl} = A\}$,
    endowed with the subspace topology induced by the hull-kernel topology on $\mathcal{P}(S)$, is a spectral space and a retrocompact subset of $\mathcal{P}(S)$ if and only if $\mathrm{cl}$ is algebraic.
\end{theorem}

\begin{proof}
    Assume first that $X$ is a spectral space and a retrocompact subset of $\mathcal{P}(S)$. To show that $\mathrm{cl}$ is algebraic, fix $A \subseteq S$. The inclusion $\bigcup\{F^\mathrm{cl} \mid F \in \mathcal{P}_\mathrm{fin}(A)\} \subseteq A^\mathrm{cl}$ follows from order-preservation, so it suffices to prove the reverse inclusion. Let $x\in A^\mathrm{cl}$. Then $V_X(A) \subseteq V_X(x)$ and thus $D_X(x) \subseteq D_X(A)$. It is easy to observe that $D_X(A) = \bigcup_{a \in A}D_X(a)$, so $\{D_X(a)\}_{a \in A}$ is an open cover of $D_X(x)$. Now $D(x)$ is a quasi-compact open subset of $\mathcal{P}(S)$ and $D_X(x) = D(x) \cap X$ is also quasi-compact, since $X$ is retrocompact in $\mathcal{P}(S)$. Hence there exists a finite subset $F \subseteq A$ such that 
    \begin{equation*}
        D_X(x) \subseteq \bigcup_{a \in F} D_X(a) = D_X(F), 
    \end{equation*}
    which yields $V_X(F) \subseteq V_X(x)$. Since $F^\mathrm{cl} \in X$ and $F \subseteq F^\mathrm{cl}$, we have $F^\mathrm{cl} \in V_X(F) \subseteq V_X(x)$, so $x \in F^\mathrm{cl}$. Therefore, every $x \in A^\mathrm{cl}$ belongs to $F^\mathrm{cl}$ for some finite subset $F \subseteq A$ and hence $A^\mathrm{cl} \subseteq \bigcup\{F^\mathrm{cl} \mid F \in \mathcal{P}_\mathrm{fin}(A)\}$. Thus, $\mathrm{cl}$ is algebraic. 
    
    Conversely, assume that $\mathrm{cl}$ is algebraic. Since $\mathcal{P}(S)$ with the hull-kernel topology is $T_0$ and $X$ carries the subspace topology, the space $X$ is also $T_0$. Now we show that $X$ is spectral by applying the ultrafilter criterion (Theorem \ref{ultrafilter_criterion}). Consider the subbasis $$\mathbb{S} := \{D_X(F) \mid F \in\mathcal{P}_\mathrm{fin}(S)\}.$$ Let $\mathscr{U}$ be an ultrafilter on $X$ and define $$A_\mathscr{U} := \{x \in S \mid V_X(x) \in \mathscr{U}\}.$$ We claim that $A_\mathscr{U} \in X_\mathbb{S}(\mathscr{U})$. We first show that $A_\mathscr{U}$ is $\mathrm{cl}$-closed. Let $x \in A_\mathscr{U}^\mathrm{cl}$. Since $\mathrm{cl}$ is algebraic, there exists a finite subset $F \subseteq A_\mathscr{U}$ such that $x \in F^\mathrm{cl}$. By definition of $A_\mathscr{U}$, $V_X(f) \in \mathscr{U}$ for every $f \in F$ and hence $$V_X(F) = \bigcap_{f \in F} V_X(f) \in \mathscr{U}.$$ Since every $\mathrm{cl}$-closed set containing $F$ also contains $x$, we have $V_X(F) \subseteq V_X(x)$. Hence $V_X(x) \in \mathscr{U}$, so $x \in A_\mathscr{U}$ and in consequence $A_\mathscr{U}^\mathrm{cl} \subseteq A_\mathscr{U}$. The reverse inclusion follows from the extension property of $\mathrm{cl}$. Thus, $A_\mathscr{U} \in X$. 
    
    Now let $F$ be a finite subset of $S$. Since $D_X(F) = \bigcup_{f \in F} D_X(f)$ and $\mathscr{U}$ is an ultrafilter, we have
    \begin{align*} 
        A_\mathscr{U} \in D_X(F) &\iff \exists f \in F \text{ such that } f \not\in A_\mathscr{U} \\ &\iff \exists f \in F \text{ such that } V_X(f) \not\in\mathscr{U} \\ &\iff \exists f \in F \text{ such that } D_X(f) \in \mathscr{U} \\ &\iff  
        D_X(F) \in \mathscr{U}. 
    \end{align*} 
    Hence $A_\mathscr{U}\in X_\mathbb{S}(\mathscr{U})$, and therefore $X$ is spectral.
    
    It remains to prove that $X$ is a retrocompact subset of $\mathcal{P}(S)$. Now let $U$ be a quasi-compact open subset of $\mathcal{P}(S)$. Then $$U = \bigcup_{k = 1}^n \bigcap_{l = 1}^{m_k} D(F_{kl})$$ for some finite sets $F_{kl} \subseteq S$. We get 
    \begin{equation*}
        U \cap X = \bigcup_{k = 1}^n \bigcap_{l = 1}^{m_k} \big(D(F_{kl}) \cap X\big) = \bigcup_{k = 1}^n\bigcap_{l = 1}^{m_k} D_X(F_{kl}). 
    \end{equation*}
    By Lemma \ref{quasi_compactness_of_subbasic_opens}, each $D_X(F_{kl})$ is quasi-compact in $X$. Since $X$ is spectral, finite intersections of quasi-compact open subsets of $X$ are again quasi-compact and open. Hence for every $k \in \{1, \ldots, n\}$ the set $\bigcap_{l = 1}^{m_k} D_X(F_{kl})$ is quasi-compact, and therefore $U \cap X$ is quasi-compact as a finite union of quasi-compact subsets. Thus, $X$ is a retrocompact subset of $\mathcal{P}(S)$. 
\end{proof}

Our next example shows that the algebraicity assumption in Theorem \ref{construction_of_spectral_space} is essential: if the closure operator is not algebraic, then the corresponding space $X$ may not be spectral. 
\begin{example}
    Let $S$ be the real unit interval, i.e. $S = [0,1]$, with the semigroup operations $x\cdot y = x \lor y := \min(x,y)$. A map $\mathrm{cl} : \mathcal{P}(S) \rightarrow \mathcal{P}(S)$ defined by
    \begin{equation*}
    A^{\mathrm{cl}} := 
    \begin{cases}
        \varnothing, & \text{for } A = \varnothing, \\
        [0, \sup(A)], & \text{for } A \neq \varnothing
    \end{cases}
    \end{equation*}
    is a closure operator on $S$. We claim that it is not algebraic. Indeed, if we take \begin{equation*}
        A := \bigg\{1 - \frac{1}{n} \hspace{3pt} \bigg| \hspace{3pt} n \in \mathbb{N}_{+}\bigg\},
    \end{equation*}
    then $\sup(A) = 1$, so $A^{\mathrm{cl}} = [0,1]$. On the other hand, 
    \begin{equation*}
        \bigcup\{F^{\mathrm{cl}} \mid F \in \mathcal{P}_\mathrm{fin}(A)\} = \bigcup_{n \in\mathbb{N}_+} \bigg[0, 1 - \frac{1}{n}\bigg] = [0, 1), 
    \end{equation*}
    hence $\mathrm{cl}$ fails algebraicity. 
    
    Let $X$ be defined as in Theorem \ref{construction_of_spectral_space} and let $F \neq \varnothing$ be a finite subset of $S$. Then \begin{equation*}
        X = \{\varnothing\} \cup \{[0,x] \mid x \in [0,1]\}
    \end{equation*}
    and 
    \begin{equation*}
        D_X(F) = \{\varnothing\} \cup \{[0, x] \mid 0 \leqslant x < \sup(F)\}. 
    \end{equation*}
    Hence an open subbasis of $X$ is given by the sets of the form $$D_t := \{\varnothing\} \cup \{[0,x] \mid 0 \leqslant x < t\}$$ for $t \in [0, 1].$ We claim that every nonempty proper open subset of $X$ is of the form $D_t$ for some $t\in[0,1]$. Indeed, let $U \subsetneq X$ be an open set. Since the family $\{D_t \mid t \in[0,1]\}$ is a subbasis, $U$ is a union of finite intersections of sets of the form $D_t$. But for all $s, t \in [0,1]$ we have $D_s \cap D_t = D_{\min(s, t)}$, and therefore every finite intersection of such sets is again of the form $D_t$. Hence for some family $\{t_i\}_{i \in I} \subseteq [0,1]$ we have $$U = \bigcup_{i\in I} D_{t_i} = D_{\sup_{i \in I}t_i}.$$
    
    Let us recall that a spectral space must have a basis of quasi-compact open sets. But for every $t > 0$ the set $D_t$ is not quasi-compact. Indeed, let $t_n$ be an increasing sequence such that $t_n \rightarrow t$ and $t_n < t$ for every natural number $n$. Then $D_t = \bigcup_{n \in \mathbb{N}} D_{t_n}$ and no finite subfamily covers $D_t$ because if $n_1, \ldots, n_k$ are fixed, then 
    \begin{equation*}
        D_{t_{n_1}} \cup \ldots \cup D_{t_{n_k}} = D_{\max(t_{n_1}, \ldots, t_{n_k})} \subsetneq D_t. 
    \end{equation*}
    Thus, $X$ is not spectral. 
\end{example}

\begin{thebibliography}{99}
\bibitem[]{Banerjee} A. Banerjee. Closure operators in abelian categories and spectral spaces. Theory and Applications of Categories 32 (2017) 719-735
\bibitem[]{Finocchiaro03042014} C. A. Finocchiaro. Spectral Spaces and Ultrafilters. Communications in Algebra 42 (2014) 1496-1508
\bibitem[]{FinocchiaroFontanaSpirito} C. A. Finocchiaro; M. Fontana; D. Spirito. Spectral spaces of semistar operations. Journal of Pure and Applied Algebra 220 (2016) 2897-2913
\bibitem[]{FinocchiaroFontanaSpiritoNullstellensatz} C. A. Finocchiaro; M. Fontana; D. Spirito. A topological version of {H}ilbert's {N}ullstellensatz. Journal of Algebra 461 (2016) 25-41
\bibitem[]{JunRayTolliver} J. Jun; S. Ray; J. Tolliver. Lattices, Spectral Spaces, and Closure Operations on Idempotent Semirings. Journal of Algebra 594 (2022) 313-363
\end{thebibliography}
\end{document}
